% GENERATED FILE. Edit canonical lessons, then run npm run generate:books.
% canonical-source-hash: fnv1a64:b3fac7ab3ce358bd
% canonical-lessons: MR-C271-kaceri, MR-C271-basasthanak, MR-C271-pangharun, MR-C271-osri, MR-C271-jhadu

\chapter{Government office, Bus station, Blanket, Veranda, Broom}
\label{ch:mr-government-office-bus-station-blanket-veranda-broom}
\hlchaptermodality{271}

\begin{chapteropening}
\textbf{By the end of this chapter:} \emph{I can say a government office, a bus station, a blanket, a veranda and a broom.}

Complete the last lesson of chapter 271: I can say a government office, a bus station, a blanket, a veranda and a broom.
\end{chapteropening}

\section[\texorpdfstring{\mr{कचेरी} (kacerī)}{कचेरी (kacerī)}]{\mr{कचेरी} (kacerī) — a government office}

\label{lesson:MR-C271-kaceri}



\begin{quote}
Before the new one: say the Marathi for a dhoti, then the Marathi for a kurta.
\end{quote}

\subsection*{You'll want to know: \mr{कचेरी}}
\textbf{\mr{कचेरी}} — \emph{kacerī} — \textquotedblleft{}a government office\textquotedblright{}.

\begin{grammarlens}[title={What you've built}]
One more word for this chapter.
\end{grammarlens}

\subsection*{Your turn}
\begin{itemize}
\raggedright
  \item \emph{Say it:} \emph{kacerī}
  \item \emph{Say it:} \emph{kacerī}, once more
  \item \emph{Say it:} say \emph{kacerī}
  \item \emph{Recall:} say the Marathi for a dhoti, then the Marathi for a kurta, then say \emph{kacerī} again
\end{itemize}

\begin{tcolorbox}[breakable,colback=teal!4,colframe=teal!35!black,title={Before you move on}]
What does \mr{कचेरी} mean? (\textquotedblleft{}A government office\textquotedblright{}.) Say it once more.
\end{tcolorbox}
\section[\texorpdfstring{\mr{बसस्थानक} (basasthānak)}{बसस्थानक (basasthānak)}]{\mr{बसस्थानक} (basasthānak) — a bus station}

\label{lesson:MR-C271-basasthanak}



\begin{quote}
Before the new one: say the Marathi for a kurta, then the Marathi for a government office.
\end{quote}

\subsection*{You'll want to know: \mr{बसस्थानक}}
\textbf{\mr{बसस्थानक}} — \emph{basasthānak} — \textquotedblleft{}a bus station\textquotedblright{}.

\begin{grammarlens}[title={What you've built}]
One more word for this chapter.
\end{grammarlens}

\subsection*{Your turn}
\begin{itemize}
\raggedright
  \item \emph{Say it:} \emph{basasthānak}
  \item \emph{Say it:} \emph{basasthānak}, once more
  \item \emph{Say it:} say \emph{basasthānak}
  \item \emph{Recall:} say the Marathi for a kurta, then the Marathi for a government office, then say \emph{basasthānak} again
\end{itemize}

\begin{tcolorbox}[breakable,colback=teal!4,colframe=teal!35!black,title={Before you move on}]
What does \mr{बसस्थानक} mean? (\textquotedblleft{}A bus station\textquotedblright{}.) Say it once more.
\end{tcolorbox}
\section[\texorpdfstring{\mr{पांघरूण} (pāṅgharūṇ)}{पांघरूण (pāṅgharūṇ)}]{\mr{पांघरूण} (pāṅgharūṇ) — a blanket}

\label{lesson:MR-C271-pangharun}



\begin{quote}
Before the new one: say the Marathi for a government office, then the Marathi for a bus station.
\end{quote}

\subsection*{You'll want to know: \mr{पांघरूण}}
\textbf{\mr{पांघरूण}} — \emph{pāṅgharūṇ} — \textquotedblleft{}a blanket\textquotedblright{}.

\begin{grammarlens}[title={What you've built}]
One more word for this chapter.
\end{grammarlens}

\subsection*{Your turn}
\begin{itemize}
\raggedright
  \item \emph{Say it:} \emph{pāṅgharūṇ}
  \item \emph{Say it:} \emph{pāṅgharūṇ}, once more
  \item \emph{Say it:} say \emph{pāṅgharūṇ}
  \item \emph{Recall:} say the Marathi for a government office, then the Marathi for a bus station, then say \emph{pāṅgharūṇ} again
\end{itemize}

\begin{tcolorbox}[breakable,colback=teal!4,colframe=teal!35!black,title={Before you move on}]
What does \mr{पांघरूण} mean? (\textquotedblleft{}A blanket\textquotedblright{}.) Say it once more.
\end{tcolorbox}
\section[\texorpdfstring{\mr{ओसरी} (osrī)}{ओसरी (osrī)}]{\mr{ओसरी} (osrī) — a veranda}

\label{lesson:MR-C271-osri}



\begin{quote}
Before the new one: say the Marathi for a bus station, then the Marathi for a blanket.
\end{quote}

\subsection*{You'll want to know: \mr{ओसरी}}
\textbf{\mr{ओसरी}} — \emph{osrī} — \textquotedblleft{}a veranda\textquotedblright{}.

\begin{grammarlens}[title={What you've built}]
One more word for this chapter.
\end{grammarlens}

\subsection*{Your turn}
\begin{itemize}
\raggedright
  \item \emph{Say it:} \emph{osrī}
  \item \emph{Say it:} \emph{osrī}, once more
  \item \emph{Say it:} say \emph{osrī}
  \item \emph{Recall:} say the Marathi for a bus station, then the Marathi for a blanket, then say \emph{osrī} again
\end{itemize}

\begin{tcolorbox}[breakable,colback=teal!4,colframe=teal!35!black,title={Before you move on}]
What does \mr{ओसरी} mean? (\textquotedblleft{}A veranda\textquotedblright{}.) Say it once more.
\end{tcolorbox}
\section[\texorpdfstring{\mr{झाडू} (jhāḍū)}{झाडू (jhāḍū)}]{\mr{झाडू} (jhāḍū) — a broom}

\label{lesson:MR-C271-jhadu}



\begin{quote}
Before the new one: say the Marathi for a blanket, then the Marathi for a veranda.
\end{quote}

\subsection*{You'll want to know: \mr{झाडू}}
\textbf{\mr{झाडू}} — \emph{jhāḍū} — \textquotedblleft{}a broom\textquotedblright{}.

\begin{grammarlens}[title={What you've built}]
One more word for this chapter.
\end{grammarlens}

\subsection*{Your turn}
\begin{itemize}
\raggedright
  \item \emph{Say it:} \emph{jhāḍū}
  \item \emph{Say it:} \emph{jhāḍū}, once more
  \item \emph{Say it:} say \emph{jhāḍū}
  \item \emph{Recall:} say the Marathi for a blanket, then the Marathi for a veranda, then say \emph{jhāḍū} again
\end{itemize}

\begin{tcolorbox}[breakable,colback=teal!4,colframe=teal!35!black,title={Before you move on}]
What does \mr{झाडू} mean? (\textquotedblleft{}A broom\textquotedblright{}.) Say it once more.
\end{tcolorbox}
